\subsection{系统构成与实现路径}
从软硬件系统构成看，知识驱动的人机协同动力电池包智能装配与在线检测模式可分为硬件层、数据层、知识层和应用层四个部分。硬件层包括工业机器人、协作机器人、视觉相机、焊接与紧固设备、力/位移传感器、条码或 RFID 采集装置、边缘控制器以及 MES/PLC 接口；其中，机器人、视觉工位与传感器单元共同承担现场执行与感知任务。软件侧主要由 MES 配方管理、质量数据库、规则推理引擎、知识图谱服务和大模型智能体服务构成，用于完成产品追溯、知识组织、异常分析和人机交互。知识层由工艺知识图谱、质量规则库和设备能力模型构成，用以表达“某类零件属于哪个装配层级”“某工步前置条件是什么”“某缺陷可能关联哪些原因”。应用层则提供工艺问答、异常诊断、换型辅助、调度建议、返修路径推荐和人机协同工位操作辅助等功能 \cite{Sahota2025JIM,Elbasheer2025JIM,Zhang2023JMSYKG,Yao2023JMSY,Long2025IJPR}。

其基本数据流与控制流可以描述为：产品进入工位后，系统首先绑定产品 ID 和工艺配方；机器人与人工依据工艺计划执行装配；视觉和传感器实时采集位置、缺陷、焊缝、力矩和状态信息；采集结果一方面进入质量数据库，另一方面通过知识映射写入知识图谱上下文；当系统发现偏差或异常时，规则引擎和大模型智能体联合给出原因解释、风险等级和处置建议；最终由操作员或上位调度系统决定继续装配、返修或重检。由此，系统将“检测发现问题”进一步扩展为“检测触发知识推理和工艺调整”的闭环机制 \cite{Sgarbossa2024JMSY,Ji2024SciRep}。

实现这一模式的关键技术主要包括以下五类。第一，装配过程知识建模技术，需要将 BOM、工艺卡、设备能力、质量规则和历史缺陷案例统一映射到可计算的知识图谱中。第二，工位级多源感知融合技术，需要将视觉、力觉、位姿和设备参数转换为可供质量判定和异常解释使用的特征。第三，任务分类与协作调度技术，需要根据任务复杂度、风险等级和人体工效要求合理划分人机分工 \cite{Faccio2021IJAMT,Alessio2024JIM,Banziger2023RCIM}。第四，面向制造场景的大模型提示与工具调用技术，需要保证大模型输出受工艺规则约束，而不是自由生成不可靠建议。第五，多智能体规划与安全控制技术，需要确保建议层的灵活性不会破坏执行层的稳定性 \cite{Massouh2025JIM,Pulikottil2023IJAMT,Abelha2020RAL}。这五类技术并不是彼此独立堆叠的模块，而是共同构成“感知--表达--推理--决策--执行”的系统链路。

若从企业导入路径考虑，该模式不宜一次性整体替换现有产线，而应分阶段实施。第一阶段可先完成工位数据采集和追溯打通；第二阶段建立缺陷库、工艺规则库和知识图谱；第三阶段在局部工位引入视觉检测和协作装配；第四阶段再上线大模型智能体，用于知识问答、工艺辅助和异常诊断建议。这种渐进式路线成本更可控、导入风险更低，也更符合制造企业引入新技术的实际节奏 \cite{Ayerbe2021AEM,Tao2020Sustainability}。从实现途径上看，本文提出的新模式并不是“另起炉灶”，而是对现有数字化装配单元的渐进式增强。

\begin{figure}[H]
  \centering
  \includegraphics[width=0.78\textwidth]{fig12_nl_simulation_framework.png}
  \caption{自然语言驱动的制造仿真框架}
  \label{fig:nl-simulation-framework}
\end{figure}

除知识组织与现场执行外，系统实现还应补上仿真验证这一中间环节。对于动力电池包装配而言，自然语言驱动的仿真框架并不是为了替代 MES 或 PLC，而是提供一种更高层的方案验证能力：工程人员先用自然语言描述资源、工序和约束，再由知识抽取模块将其转化为结构化参数，进一步生成可执行仿真模型并输出节拍、瓶颈和资源利用结果。图~\ref{fig:nl-simulation-framework}展示了这一从用户请求到仿真执行的四阶段链路，说明大模型辅助不只可以停留在问答层，还可以延伸到生产组织、换型评估和方案预验证 \cite{Elbasheer2025JIM,Long2025IJPR}。

\begin{figure}[H]
  \centering
  \includegraphics[width=0.7\textwidth]{fig5_kg_construction.png}
  \caption{装配文本知识图谱构建流程}
  \label{fig:mode-placeholder}
\end{figure}

图~\ref{fig:mode-placeholder}并不直接表示整条产线的软硬件系统架构，而是用于说明装配知识图谱的构建逻辑。图中可以看到本体构建、实体识别、关系抽取、知识融合以及数据库生成等关键步骤，说明知识层并不是凭空得到的，而是要把装配文档、工艺描述和历史数据转化为可计算、可查询的结构化知识。因而，该图可作为本文“知识层如何形成”的具体支撑 \cite{Wang2025CJME}。
